%! Algebra 2 — Mid-Year Final Exam — Study Guide (reference sheet, no problems, no key)
%
% Course-level, not unit-level: the final covers Units 1–2 and Lessons 3.1–3.2, so it lives in
% final/ (built by final/Makefile), not in a unit's study_guide/ — unit.mk would merge a unit
% study guide into that unit's packets.
%
% Same shape as the unit study guides (unit01/study_guide is the reference): a REFERENCE sheet —
% vocabulary, forms, and each lesson's target misconception as a caution. No problems, so no
% key: the problem set is the practice final, a printed copy of the DeltaMath test
% "A2: Mid-Year Final Exam" (47 questions, same order). The exam map below gives the question
% numbers per lesson; keep it in step with COURSE_BREAKDOWN.md if the DeltaMath test changes.
%
% Not on the final, so not here: completing the square (user, 2026-11-27), and Lessons 3.3–3.6.
%
% No \namedateperiod: neither a lesson cover nor a test. FOUR PAGES at 10pt is the ceiling
% (two sheets, duplex).
\documentclass[10pt]{article}
\usepackage{algebra2-article}
\usepackage{algebra2-boxes}

% ── Local layout macros (as in the unit study guides) ────────────────────────
\newcommand{\sghead}[2]{%
  \boxguard[10]%
  \vspace{4pt}\par\noindent
  \colorbox{forest}{\parbox{\dimexpr\linewidth-2\fboxsep}{%
    \small\color{white}\bfseries #1 \hfill \textmd{\itshape\scriptsize\color{forestmid} #2}}}%
  \vspace{2pt}\par}

\newcommand{\sgwatch}[1]{%
  \vspace{1pt}\par\noindent
  \colorbox{redbg}{\parbox{\dimexpr\linewidth-2\fboxsep}{%
    \footnotesize\textbf{\textcolor{redacc}{Watch out:}} #1}}%
  \vspace{1pt}\par}

% tabularx must see a literal \end{tabularx}, so the spec is a column type, not an environment.
\newcolumntype{G}{>{\bfseries\leavevmode\color{forest}}p{0.245\linewidth}}
\newcommand{\sgtabsetup}{\renewcommand{\arraystretch}{1.06}\footnotesize}

\begin{document}

\pageheader{Mid-Year Final Exam}{Study Guide --- Reference Sheet}

\vspace{-0.30in}
\noindent\footnotesize The final is a \textbf{DeltaMath test}: \textbf{47 questions},
\textbf{90 minutes}, calculator allowed. It covers \textbf{Units 1 and 2} and \textbf{Lessons 3.1
and 3.2}. The \textbf{practice final} has the same questions with new numbers, in the same order.
Work it first. Then use the map to find the lesson behind any question you missed, and read that
lesson's section below. There are no problems on this sheet. It tells you what to \emph{know}.
\par\normalsize

% ======================================================================
\sghead{The exam map}{question numbers on the final and the practice final}
{\sgtabsetup
\begin{tabularx}{\linewidth}{@{} >{\bfseries\leavevmode\color{forest}}p{0.11\linewidth} >{\bfseries}p{0.30\linewidth} X @{}}
1--2   & 1.1 Lines & Equation from a point and slope; equation from two points on a graph. \\
3--5   & 1.2 Absolute value equations \& inequalities & Solve $|X|=c$ (two levels); solve $|X|<c$ or $|X|>c$. \\
6      & 1.3 Absolute value functions & Read vertex form $a|x-h|+k$. \\
7--9   & 1.4 Piecewise \& step functions & Evaluate from the rule and from a graph; read a step function. \\
10     & 1.5 Linear regression & Line of best fit and $r$ from data (calculator). \\
11--15 & 2.0--2.1 Quadratic features \& graphs & Features; vertex and intercept form; vertex form to standard; graph and state properties. \\
16--20 & 2.2 Factoring \& solving & Factor $ax^2+bx+c$ ($a>1$); solve by factoring; solve $x^2=kx$-type equations. \\
21--24 & 2.3 Square roots \& $i$ & Solve with square roots; simplify $\sqrt{-k}$; add and multiply imaginary numbers. \\
25--27 & 2.4 Complex numbers & Add, subtract, multiply, and square $a+bi$. \\
28--30 & Quadratic formula \& discriminant & Real roots; complex roots; number and type of roots. \\
31--35 & 2.5 Systems \& modeling & Line--parabola systems; read a quadratic in context; word problems. \\
36--40 & 3.1 Polynomial operations & Add/subtract; square a binomial; multiply; factored to standard form. \\
41--47 & 3.2 Factoring polynomials & GCF; grouping; sum and difference of cubes; quartic trinomials; factoring completely. \\
\end{tabularx}}

% ======================================================================
\sghead{1.1 \quad Lines: Slope, Forms \& Graphing}{A.F.1a--e \;\textbullet\; Questions 1--2}
{\sgtabsetup
\begin{tabularx}{\linewidth}{@{} G X @{}}
Slope & $m=\dfrac{y_2-y_1}{x_2-x_1}$, the rise over the run. Any two points on the line give the
  same answer. In context it is a rate, with units. \\
Slope-intercept & $y=mx+b$: the rate $m$ and the starting value $b$. \\
Point-slope & $y-y_1=m(x-x_1)$ --- from a slope and \emph{any} point. Expand it to reach
  $y=mx+b$. \\
Reading a graph & Pick two lattice points (where the line crosses grid corners), count rise over
  run, then read or solve for $b$. \\
Parallel / perpendicular & Equal slopes; opposite reciprocal slopes ($m_1m_2=-1$). \\
Special cases & Horizontal $y=c$ ($m=0$); vertical $x=c$ (slope undefined). \\
\end{tabularx}}
\sgwatch{Keep the subtraction in the same order top and bottom. $\frac{y_2-y_1}{x_1-x_2}$ flips
the sign of the slope. And $y-52=3(x-4)$ and $y=3x+40$ are the \emph{same} line --- expand
before you decide two answers differ.}

% ======================================================================
\sghead{1.2 \quad Absolute Value Equations \& Inequalities}{A2.EI.1a--e \;\textbullet\; Questions 3--5}
{\sgtabsetup
\begin{tabularx}{\linewidth}{@{} G X @{}}
Isolate first & Get the bars alone before you split: $3|x+1|-2=10$ becomes $|x+1|=4$
  \emph{first}. \\
$|X|=c$ & $X=c$ \textbf{or} $X=-c$. Two solutions if $c>0$, one if $c=0$, \textbf{none} if
  $c<0$. \\
$|X|<c$ & ``Less th\textbf{AND}'': one interval, $-c<X<c$. \\
$|X|>c$ & ``Great\textbf{OR}'': two rays, $X<-c$ \textbf{or} $X>c$. \\
Notation & Interval $(-4,4)$ or $[-4,4]$; a number line with open or closed dots. $\infty$
  never takes a bracket. Dividing by a negative flips the inequality. \\
\end{tabularx}}
\sgwatch{Splitting \emph{before} isolating turns $3|x+1|-2=10$ into $3x+1-2=10$ --- one answer,
and a wrong one. And $<$ where you needed $>$ gives the \emph{opposite} set, not a slightly
different one.}

% ======================================================================
\sghead{1.3 \quad Absolute Value Functions \& Transformations}{A2.F.1b--c; A2.F.2 \;\textbullet\; Question 6}
{\sgtabsetup
\begin{tabularx}{\linewidth}{@{} G X @{}}
Vertex form & $g(x)=a|x-h|+k$: vertex $(h,k)$, axis of symmetry $x=h$. \\
Finding $h$ & Set the inside of the bars to zero. Inside is horizontal and backwards: $|x-3|$
  moves \emph{right} $3$; $|x+3|$ moves left $3$. \\
Sign of $a$ & $a>0$ opens up and $k$ is a \textbf{minimum}; $a<0$ opens down and $k$ is a
  \textbf{maximum}. \\
Size of $a$ & $|a|$ is the slope of the right-hand ray: $|a|>1$ narrower, $0<|a|<1$ wider. \\
Range & $y\ge k$ (opens up) or $y\le k$ (opens down). Domain is all reals. \\
\end{tabularx}}
\sgwatch{A negative $a$ changes the vertex's \emph{job}. For $R(x)=-2|x-3|+7$, $7$ is the
\emph{maximum} --- check: $R(5)=3$, which is less than $7$.}

% ======================================================================
\sghead{1.4 \quad Piecewise \& Step Functions}{A2.F.2a--c, f \;\textbullet\; Questions 7--9}
{\sgtabsetup
\begin{tabularx}{\linewidth}{@{} G X @{}}
Evaluating & First find the piece whose condition the input satisfies, \emph{then} use that rule.
  Never use both. \\
At a boundary & Exactly one piece owns the input --- the one whose inequality \textbf{includes} it
  ($\le$ or $\ge$). On a graph, that is the \textbf{closed} dot. \\
From a graph & To read $f(2)$, go to $x=2$ and find the closed (or unbroken) point above or below
  it; an open dot is \emph{not} the value. \\
Step function & Constant on each interval, then jumps. $\lfloor x\rfloor$ rounds \emph{down}:
  $\lfloor 3.9\rfloor=3$. A ``started hour'' charge rounds \emph{up}. \\
\end{tabularx}}
\sgwatch{Rounding down is not dropping the decimal: $\lfloor -2.5\rfloor=-3$, not $-2$. On the
negative side, down means \emph{more} negative.}

% ======================================================================
\sghead{1.5 \quad Linear Regression}{A2.ST.2c--h \;\textbullet\; Question 10}
{\sgtabsetup
\begin{tabularx}{\linewidth}{@{} G X @{}}
On the calculator & Enter $x$ in L1 and $y$ in L2, then run \textbf{LinReg($ax+b$)}. Round as the
  question asks. \\
Line of best fit & $\hat y = ax+b$; the hat means \emph{predicted}. The slope is a per-unit rate,
  the intercept the value at $x=0$ --- both with units. \\
Correlation $r$ & Between $-1$ and $1$: the sign is the direction, $|r|$ the strength. \\
Predictions & Inside the data's range (interpolation) is reasonable; outside it (extrapolation)
  may break. A strong $r$ never proves cause. \\
\end{tabularx}}
\sgwatch{$r=-0.9$ is a \emph{stronger} fit than $r=0.4$. Compare $|r|$, never $r$ itself.}

% ======================================================================
\sghead{2.0--2.1 \quad Quadratic Features, Forms \& Graphs}{A2.F.1; A2.F.2 \;\textbullet\; Questions 11--15}
{\sgtabsetup
\begin{tabularx}{\linewidth}{@{} G X @{}}
Vertex & The turning point $(h,k)$: a \textbf{minimum} if it opens up ($a>0$), a
  \textbf{maximum} if it opens down ($a<0$). \\
Axis of symmetry & The line $x=h$. Points equally far either side have equal outputs. \\
Vertex form & $y=a(x-h)^2+k$: vertex $(h,k)$. $(x+2)^2$ is $(x-(-2))^2$, so $h=-2$. \\
Standard form & $y=ax^2+bx+c$: $y$-intercept $(0,c)$; axis $x=-\frac{b}{2a}$; put that $x$ back in
  for $k$. \\
Intercept form & $y=a(x-p)(x-q)$: zeros $p$ and $q$; axis halfway, $x=\frac{p+q}{2}$. \\
Vertex $\to$ standard & Square the binomial, distribute $a$, add $k$:
  $2(x-1)^2+3=2(x^2-2x+1)+3=2x^2-4x+5$. \\
Range, ends, intervals & Range $y\ge k$ or $y\le k$. Both ends go the same way. Increasing and
  decreasing split \emph{at} $x=h$. \\
\end{tabularx}}
\sgwatch{The sign in the bracket flips. $(x-1)^2$ moves \emph{right} $1$, and $(x+1)(x-3)$ has
zeros $-1$ and $3$. Set each bracket to zero; never copy the printed sign.}

% ======================================================================
\sghead{2.2 \quad Factoring \& Solving by Factoring}{A2.EO.3b, d; A2.EI.2b, d \;\textbullet\; Questions 16--20}
{\sgtabsetup
\begin{tabularx}{\linewidth}{@{} G X @{}}
GCF first & Always: $6x^2+9x=3x(2x+3)$. \\
$a>1$ trinomials & For $ax^2+bx+c$, find two numbers with product $ac$ and sum $b$, split $bx$,
  and group: $2x^2+7x+3=2x^2+6x+x+3=(2x+1)(x+3)$. Check by multiplying back. \\
Special patterns & $x^2-k^2=(x+k)(x-k)$; $x^2+2kx+k^2=(x+k)^2$. A sum of squares $x^2+k^2$ does not
  factor over the reals. \\
Zero Product Property & If $AB=0$, then $A=0$ or $B=0$. So: set $=0$, factor, set each factor to
  zero. \\
\end{tabularx}}
\sgwatch{The property needs \emph{zero}. Never divide $x^2=5x$ by $x$ --- that loses $x=0$. Write
$x^2-5x=0$, $x(x-5)=0$, so $x=0$ or $5$. And $(x-1)(x+2)=4$ does \emph{not} mean each factor is
$4$.}

% ======================================================================
\sghead{2.3 \quad Square Roots \& the Imaginary Unit}{A2.EI.2b, d; A2.EO.4a--b \;\textbullet\; Questions 21--24}
{\sgtabsetup
\begin{tabularx}{\linewidth}{@{} G X @{}}
Square Root Property & Isolate the square, then $x^2=k$ gives $x=\pm\sqrt k$ and $(x-h)^2=k$ gives
  $x=h\pm\sqrt k$. Keep the $\pm$. \\
Simplest radical form & No perfect-square factor left inside: $\sqrt{50}=5\sqrt2$. \\
Imaginary unit & $i=\sqrt{-1}$, $i^2=-1$. For $k>0$, $\sqrt{-k}=i\sqrt k$:
  $\sqrt{-18}=3i\sqrt2$. \\
Adding imaginaries & Like terms: $\sqrt{-9}+\sqrt{-16}=3i+4i=7i$. \\
\end{tabularx}}
\sgwatch{Pull the $i$ out of each negative radical \emph{before} you multiply:
$\sqrt{-4}\cdot\sqrt{-9}=2i\cdot3i=6i^2=-6$, not $\sqrt{36}=6$.}

% ======================================================================
\sghead{2.4 \quad Operations with Complex Numbers}{A2.EO.4c \;\textbullet\; Questions 25--27}
{\sgtabsetup
\begin{tabularx}{\linewidth}{@{} G X @{}}
Complex number & $a+bi$ in standard form: real part $a$, imaginary part $b$. \\
Add / subtract & Real with real, imaginary with imaginary. Distribute the minus to \emph{both}
  parts: $(5+2i)-(3-4i)=2+6i$. \\
Multiply & Distribute, then replace $i^2$ with $-1$: $(2+i)(3-i)=6+i-i^2=7+i$. \\
Squaring & $(a+bi)^2=(a+bi)(a+bi)$: $(3+2i)^2=9+12i+4i^2=5+12i$. \\
Conjugates & $(a+bi)(a-bi)=a^2+b^2$, always real. \\
\end{tabularx}}
\sgwatch{$(3+2i)^2$ is not $9+4i^2$ --- the middle term $12i$ comes from the cross products. Write
the square out as two factors. And no $i^2$ in a final answer.}

% ======================================================================
\sghead{The Quadratic Formula \& the Discriminant}{A2.EI.2b, d; A2.F.2d \;\textbullet\; Questions 28--30}
{\sgtabsetup
\begin{tabularx}{\linewidth}{@{} G X @{}}
Standard form first & $ax^2+bx+c=0$ --- everything on one side; read $a$, $b$, $c$ \emph{with
  their signs}. \\
Quadratic formula & $x=\dfrac{-b\pm\sqrt{b^2-4ac}}{2a}$. It solves every quadratic. Simplify the
  radical, then divide \emph{every} term by $2a$. \\
Complex roots & A negative under the radical gives $p\pm qi$:
  $x^2+2x+5=0 \Rightarrow x=\dfrac{-2\pm\sqrt{-16}}{2}=-1\pm2i$. \\
Discriminant & $D=b^2-4ac$. $D>0$: two real roots (crosses twice). $D=0$: one repeated real root
  (touches once). $D<0$: two complex roots (misses the axis). A perfect-square $D$ means rational
  roots. \\
\end{tabularx}}
\sgwatch{Square \emph{all} of $b$: for $b=-3$, $b^2=(-3)^2=9$, not $-9$. And a negative
discriminant is not ``no solution'' --- it is two \emph{complex} solutions.}

% ======================================================================
\sghead{2.5 \quad Quadratic Systems \& Modeling}{A2.EI.3a--d; A2.EI.2a, d \;\textbullet\; Questions 31--35}
{\sgtabsetup
\begin{tabularx}{\linewidth}{@{} G X @{}}
Solution of a system & An ordered pair $(x,y)$ true in \emph{both} equations --- where the graphs
  meet. \\
Substitution & Set the $y$'s equal, move everything to one side, solve, then put each $x$ into
  the \emph{line} for its $y$. \\
How many & A line and a parabola cross ($2$), touch ($1$), or miss ($0$) --- the discriminant of
  the collapsed equation. \\
Reading a model & $h(0)$: the start. Positive zero: when it lands (reject $t<0$). Vertex input
  $-\frac{b}{2a}$: \emph{when} it peaks. Vertex output: \emph{how high}. \\
\end{tabularx}}
\sgwatch{For $h(t)=-16t^2+32t+48$, ``how high does it go?'' is $64$~ft --- not $1$ (the time),
not $48$ (the start), not $3$ (the landing).}

% ======================================================================
\sghead{3.1 \quad Polynomial Operations}{A2.EO.3a \;\textbullet\; Questions 36--40}
{\sgtabsetup
\begin{tabularx}{\linewidth}{@{} G X @{}}
Standard form \& degree & Terms in descending powers; the degree is the highest power. \\
Add / subtract & Combine like terms (same variable, same power). A minus before parentheses is a
  $-1$ that reaches \emph{every} term: $(4x^2-x)-(x^2-3x+2)=3x^2+2x-2$. \\
Multiply & Every term of one factor times every term of the other (distribution or a box), then
  combine like terms. \\
Special products & $(a+b)^2=a^2+2ab+b^2$; $(a-b)^2=a^2-2ab+b^2$; $(a+b)(a-b)=a^2-b^2$. \\
Three or more factors & Multiply two at a time: $(x-1)(x+2)(x-3)$ is $(x^2+x-2)(x-3)$. Factored
  to standard form is the same job. \\
\end{tabularx}}
\sgwatch{Exponents do not distribute over a sum: $(x+2)^2\ne x^2+4$ and $(x+2)^3\ne x^3+8$. Test
with $x=1$: $(1+2)^3=27$, but $1+8=9$.}

% ======================================================================
\sghead{3.2 \quad Factoring Polynomials}{A2.EO.3b, d \;\textbullet\; Questions 41--47}
{\sgtabsetup
\begin{tabularx}{\linewidth}{@{} G X @{}}
GCF & Numbers \emph{and} variables, lowest power of each: $12x^3y-18x^2y^2=6x^2y(2x-3y)$. \\
Grouping (four terms) & Pair, take out each pair's GCF, then the common binomial:
  $x^3+3x^2-4x-12=x^2(x+3)-4(x+3)=(x+3)(x^2-4)=(x+3)(x+2)(x-2)$. \\
Sum / difference of cubes & $a^3+b^3=(a+b)(a^2-ab+b^2)$; $a^3-b^3=(a-b)(a^2+ab+b^2)$ ---
  ``SOAP'': \textbf{S}ame sign, \textbf{O}pposite, \textbf{A}lways \textbf{P}ositive. Know the
  cubes $1,8,27,64,125$. \\
Quadratic form & $x^4-5x^2+4$ is $u^2-5u+4$ with $u=x^2$: $(x^2-1)(x^2-4)$, then keep going. \\
Factor completely & GCF, then the pattern, then check every factor again. Done only when no factor
  factors further. \\
\end{tabularx}}
\sgwatch{Stopping early: $x^4-16=(x^2+4)(x^2-4)$ is not finished --- $x^2-4=(x+2)(x-2)$. But
$x^2+4$ is prime: it is \emph{not} $(x+2)^2$.}

\end{document}
